\documentclass[12pt]{article}
\usepackage[utf8]{inputenc}
\usepackage[russian]{babel}
\usepackage{amsmath}
\usepackage{amssymb}
\usepackage{geometry}
\geometry{a4paper, margin=2cm}

\begin{document}

\section*{Задача 3.1.6f}

\textbf{Условие:} Отсортировать последовательность $a_1,\ldots,a_n$ в порядке возрастания элементов.

\texttt{in:} $n,a_1,\ldots,a_n$\\
\texttt{out:} $b_1,\ldots,b_n$, где $b_1\leq\cdots\leq b_n$

\section*{Псевдокод}
\begin{verbatim}
procedure sort(l, r: longint);
var
    i, j, x, y: longint;
begin
    i := l;
    j := r;
    x := a[(l+r) div 2];
    repeat
        while a[i] < x do
            inc(i);
        while x < a[j] do
            dec(j);
        if not (i > j) then
        begin
            y := a[i];
            a[i] := a[j];
            a[j] := y;
            inc(i);
            j := j - 1;
        end;
    until i > j;
    if l < j then
        sort(l, j);
    if i < r then
        sort(i, r);
end;

begin
    max := len(array);
    sort(1, max);
end;
\end{verbatim}

\section*{Решение}

\begin{verbatim}
Read 1
Load =0
Sub 1
JZero END
JGTZ END    ; n <= 0 ==> end
Load 1
Add =10
Store 2     ; r2 := 10 + n, end of array + 1

Load =10
Store 1     ; r1 := 10, start of array. And also r1 is a helping register

READ_ARRAY: Load 2
    Sub 1
    JZero INIT_SORT ; r1 == r2 => end read
    Read *1         ; buf[r1] <- In
    Load =1
    Add 1
    Store 1         ; r1 += 1
    Jump READ_ARRAY
INIT_SORT: Load =1
    Store 3         ; r3 := l = 1
    Load 2
    Store 4         ; r4 := &r = 10 + n
    Sub =10
    Store *4        ; r (= *r4) := n
    Load 4
    Add =2
    Store 9         ; r9 := &cs
    Load =0
    Store *9        ; cs := 0 (call from MAIN)
SORT: Load 4
    Add =1
    Store 5         ; r5 := &i = &r + 1 (next i on call stack)
    Load 3
    Store *5        ; i (= *r5) := l
    Load *4
    Store 6         ; r6 := j = last r on call stack
    Load 3
    Add *4
    Div =2          ; r0 := (l+r) // 2
    Add =9          ; r0 := &a[(l+r) // 2]
    Load *0
    Store 7         ; r7 := x = a[(l+r) // 2]
REPEAT: Jump WHILE_i
WHILE_i: Load *5
    Add =9
    Load *0
    Sub 7
    JGTZ WHILE_j
    JZero WHILE_j
    Load *5
    Add =1
    Store *5
    Jump WHILE_i
WHILE_j: Load 6
    Add =9
    Load *0
    Sub 7
    Mult =-1
    JGTZ IF_not_i_gt_j
    JZero IF_not_i_gt_j
    Load 6
    Sub =1
    Store 6
    Jump WHILE_j
IF_not_i_gt_j: Load *5
    Sub 6
    JGTZ UNTIL
    Load *5
    Add =9
    Load *0
    Store 8
    Load *5
    Add =9
    Store 1
    Load 6
    Add =9
    Load *0
    Store *1
    Load 6
    Add =9
    Store 1
    Load 8
    Store *1
    Load *5
    Add =1
    Store *5
    Load 6
    Sub =1
    Store 6
UNTIL: Load *5
    Sub 6
    JGTZ IF_l_lt_j
    Jump REPEAT
IF_l_lt_j: Load 3
    Sub 6
    JGTZ IF_i_lt_r
    JZero IF_i_lt_r
    Load 4
    Add =3
    Store 4
    Load 6
    Store *4
    Load 9
    Add =3
    Store 9
    Load =1
    Store *9
    Jump SORT
IF_i_lt_r: Load *5
    Sub *4
    JGTZ CALL_STACK_RETURN
    JZero CALL_STACK_RETURN
    Load *4
    Store 1
    Load 4
    Add =3
    Store 4
    Load 1
    Store *4
    Load *5
    Store 3
    Load 5
    Add =3
    Store 5
    Load 9
    Add =3
    Store 9
    Load =2
    Store *9
    Jump SORT
CALL_STACK_RETURN: Load *9
    JZero PRINT_INIT
    Load 4
    Sub =3
    Store 4
    Load 5
    Sub =3
    Store 5
    Load *9
    Store 1
    Load 9
    Sub =3
    Store 9
    Load 1
    Sub =1
    JZero IF_i_lt_r
    Jump CALL_STACK_RETURN
PRINT_INIT: Load =10
    Store 1
PRINT: Load 2
    Sub 1
    JZero END
    Write *1
    Load 1
    Add =1
    Store 1
    Jump PRINT
END: Halt 0
\end{verbatim}

\section*{Сложность}

Рассматриваем допустимый вход $n\geq 1$. Все команды RAM, включая чтение, запись, переходы и обслуживание стека вызовов, считаются за один шаг.

До вызова \texttt{SORT} выполняется $8n+24$ шагов:
\begin{enumerate}
    \item Проверка $n$ и подготовка адресов — $10$ шагов.
    \item Чтение массива — $8n+3$ шага.
    \item Инициализация сортировки и стека — $11$ шагов.
\end{enumerate}

После завершения сортировки подготовка вывода, вывод массива и остановка занимают $8n+6$ шагов. Следовательно, суммарные затраты вне \texttt{SORT} равны
\[
    T_{\mathrm{outside}}(n)=16n+30.
\]

Стоимость разбиения зависит от входного массива. В худшем суммарном случае для отрезка длины $m\geq 5$ опорный элемент является минимальным: разбиение оставляет правую подзадачу длины $m-1$. Само разбиение занимает $11m+41$ шаг. С учётом подготовки вызова, проверок и возврата такой уровень рекурсии добавляет
\[
    13+(11m+41)+8+18+17=11m+97
\]
шагов к стоимости оставшейся подзадачи.

Для отрезка длины $4$ максимальная суммарная стоимость достигается при разбиении с двумя обменами: получаются две подзадачи длины $2$. Разбиение занимает $85$ шагов. Каждая подзадача длины $2$ занимает $73$ шага без учёта обработки её возврата; накладные расходы на левый и правый вызовы равны соответственно $34$ и $35$ шагам. Поэтому затраты \texttt{SORT} для отрезка длины $4$, без финальной проверки завершения главного вызова, равны
\[
    13+85+8+34+35+2\cdot73=321.
\]

Для $n\geq 4$ получаем
\[
\begin{aligned}
T_{\mathrm{SORT}}(n)
&=321+\sum_{m=5}^{n}(11m+97)+2\\
&=\frac{11n^2+205n-350}{2}.
\end{aligned}
\]
Здесь последние $2$ шага — проверка выхода из главного вызова.

Для малых входов подсчёт отдельно даёт:
\[
\begin{aligned}
T_{\mathrm{SORT}}(1)&=63,\\
T_{\mathrm{SORT}}(2)&=75,\\
T_{\mathrm{SORT}}(3)&=184.
\end{aligned}
\]

Добавляя чтение входа и вывод результата, получаем точное число шагов в худшем случае:
\[
\boxed{
T_{\max}(n)=
\begin{cases}
109, & n=1,\\
137, & n=2,\\
262, & n=3,\\
\dfrac{11n^2+237n-290}{2}, & n\geq 4.
\end{cases}
}
\]

Таким образом, худшая сложность этой реализации QuickSort равна
\[
    T_{\max}(n)=\Theta(n^2).
\]
Хотя QuickSort обычно имеет ожидаемую сложность $O(n\log n)$ на случайной перестановке, выбор опорного элемента по средней позиции массива не гарантирует сбалансированных разбиений.

\end{document}